\documentclass[11pt]{article}
\usepackage[a4paper,margin=27mm,headheight=14pt]{geometry}
\usepackage[T1]{fontenc}
\usepackage[utf8]{inputenc}
\usepackage{amsmath,amsthm}
\usepackage{newpxtext,newpxmath}
\usepackage{microtype,mathtools,amsthm,booktabs,longtable,array,enumitem}
\usepackage[dvipsnames]{xcolor}
\usepackage{fancyhdr,tocloft}
\setlength{\cftbeforesecskip}{4pt}
\setlength{\cftbeforetoctitleskip}{0pt}
\setlength{\cftaftertoctitleskip}{12pt}
\usepackage[hyphens]{url}
\usepackage[colorlinks=true,linkcolor=MidnightBlue,citecolor=MidnightBlue,urlcolor=MidnightBlue]{hyperref}
\hypersetup{pdftitle={Sharp Stability Constants for Fractional Covers and Incidence Spans},pdfauthor={ChatGPT; research memorandum for ProveIt},pdfsubject={Submodularity, entropy, Ramsey density, and Lagrangian intersection counts}}
\setlength{\emergencystretch}{2em}
\setlength{\parskip}{3pt}
\setlist{nosep}
\pagestyle{fancy}
\fancyhf{}
\fancyhead[L]{\small Sharp fractional-cover stability}
\fancyhead[R]{\small Research memorandum}
\fancyfoot[C]{\thepage}
\newtheorem{theorem}{Theorem}[section]
\newtheorem{lemma}[theorem]{Lemma}
\newtheorem{proposition}[theorem]{Proposition}
\newtheorem{corollary}[theorem]{Corollary}
\theoremstyle{definition}
\newtheorem{definition}[theorem]{Definition}
\newtheorem{example}[theorem]{Example}
\newtheorem{question}[theorem]{Research question}
\theoremstyle{remark}
\newtheorem{remark}[theorem]{Remark}
\newcommand{\E}{\mathbb E}
\newcommand{\Pp}{\mathbb P}
\newcommand{\Fq}{\mathbb F_q}
\newcommand{\RR}{\mathbb R}
\newcommand{\NN}{\mathbb N}
\newcommand{\1}{\mathbf 1}
\newcommand{\rank}{\operatorname{rank}}
\newcommand{\Span}{\operatorname{span}}
\newcommand{\codim}{\operatorname{codim}}
\newcommand{\KL}{D_{\mathrm{KL}}}
\newcommand{\TV}{d_{\mathrm{TV}}}
\newcommand{\Red}{\mathfrak R}
\newcommand{\Gap}{\mathfrak G}
\newcommand{\Kopt}{\mathsf K}
\newcommand{\binomtwo}[1]{\binom{#1}{2}}
\newcommand{\src}[1]{\texttt{\detokenize{#1}}}
\newcommand{\oaicommit}{\texttt{adc7f1241b42e322a6451854ab7e4b4c146bf78a}}
\newcommand{\picommit}{\texttt{1210e9d3ac88a573f98d7ef8626ba3d2c30e9749}}

\title{\bfseries Sharp Stability Constants for Fractional Covers
and Incidence Spans\\[5pt]
\large Ramsey density, equality structure, and singular Lagrangian counts}
\author{Research memorandum prepared with ChatGPT\\for the ProveIt project}
\date{October 7, 2026 \quad\textperiodcentered\quad Version 1.0}

\begin{document}
\maketitle
\begin{abstract}
We determine the exact constant comparing the total subadditivity defect of a
submodular function with a prescribed fractional-cover gap. For a fractional
cover $(S_j,\gamma_j)$ of a finite ground set $E$, the constant is
\[
 \Kopt_\gamma=\max_{F\subseteq E,\ |F|\ge2}
 \frac{|F|-1}{\sum_{j:S_j\cap F\ne\varnothing}\gamma_j-1},
\]
provided all denominators are positive. The inequality holds for every normalized
monotone submodular function, and for every normalized submodular function when
the cover is a fractional partition. A unit-contraction argument supplies a
self-contained proof. The same constant is sharp for discrete entropy and, in
the fractional-partition case, for positive-definite determinant inequalities.
For subspaces indexed by edges of a simple uniform hypergraph, the result becomes
an optimal incidence-span inequality; in graphs its constant is precisely
$m_2(G)$. We classify equality for strictly balanced hypergraphs and identify the
exact dimension cost of repairing a family to a direct sum.
A separate exterior-algebra count controls dependent families of intersections
of independent random Lagrangians. Applied to the explicit 13-point graph in the
OpenAI Sidorenko manuscript, it replaces the nondirect residual estimate
$O_D(q^{-D+1641})$ by $O_D(q^{-D+224})$ for even $D\ge364$; the residual argument
alone already gives $O_D(q^{-7})$ at $D=310$. This does not set the dimension
threshold for the full Sidorenko construction. Exact finite certificates and a
research agenda accompany the proofs.
\end{abstract}

\vspace{4pt}
\noindent\textbf{Status.} The theorems below are proved in this manuscript.
The accompanying computations are supplementary exact checks, not a formal
verification. No new Lean compilation or independent verification of either
repository's entire breakthrough program is claimed. Priority beyond the
explicitly compared literature remains to be established.

\newpage
\begingroup
\fontsize{9.5}{11.3}\selectfont
\tableofcontents
\endgroup
\newpage

\section{Repository review and the research target}
\label{sec:review}

\subsection{Snapshots, rather than an unqualified list of breakthroughs}

This study uses the following fixed repository states:
\begin{center}
\begin{tabular}{@{}p{0.25\linewidth}p{0.68\linewidth}@{}}
\toprule
Repository & Pinned commit and observed date\\
\midrule
\src{openai/math} & \oaicommit\newline October 6, 2026\\[3pt]
\src{ProveIt} & \picommit\newline October 7, 2026\\
\bottomrule
\end{tabular}
\end{center}
The OpenAI release describes a collection of 722 manuscripts in 372 result
families and explicitly distinguishes different stages of verification
\cite{OpenAIReadme,OpenAICatalog}. Its catalog contains claims concerning
quasipolynomial bounds for arithmetic progressions, superexponential van der
Waerden numbers, sharp logarithmic exponents in off-diagonal Ramsey theory, and
counterexamples to Sidorenko's conjecture. A catalog entry is evidence of what
is claimed, not a substitute for checking the claim's proof.

The ProveIt checkpoint takes a similarly useful, more fine-grained approach.
Its Ramsey status file reports 113 checked exact companions among 120 numbered
source statements, seven open statements, a successful build of 739 reachable
modules, and an axiom audit of 4,709 public theorems. It also records repairs to
geometric interfaces and states that an improvement of the general all-length
Gowers--Szemer\'edi bound is not claimed \cite{ProveItStatus}. These are
\emph{repository-reported} verification results; the present work did not rerun
that build. The checkpoint includes vendored OpenAI mathematics modules, so the
two projects already have an explicit technical connection.

The review here is deliberately selective. It examines the current status and
several candidate directions, then reads the relevant finite-complex,
rank-layer, and singular-tail arguments of the Sidorenko manuscript in detail.
It is not an audit of all 722 manuscripts, nor of all ProveIt research reports.

\subsection{Why the selected target is not a catalog-level extrapolation}

The progression and Ramsey-number claims are natural candidates, but their
principal quantitative mechanisms would require substantial independent
verification before importing them as black boxes. The selected Sidorenko
manuscript contains a more accessible structural bottleneck: it assigns
subspaces to edges of a fixed graph and uses the inequality
\begin{equation}
 \sum_v\dim\Bigl(\sum_{e\ni v}H_e\Bigr)
 \ge 2\dim\Bigl(\sum_e H_e\Bigr),
 \label{eq:sourceweak}
\end{equation}
with an extra unit of dimension when the family is not direct
\cite[\src{sections/tail.tex}, \src{lem:tail-span}]{OpenAISidorenko}.
This lemma is elementary, transparent, and independent of the difficult
finite-field construction that surrounds it.

The question we pursue is stronger than whether the extra unit can be enlarged
for one example:
\begin{quote}
For a fixed incidence pattern, what is the \emph{optimal} comparison between
local span excess and the total dimension of linear dependence? Does the same
comparison hold for arbitrary submodular functions, and can it improve the
singular counting argument rather than merely restate its qualitative step?
\end{quote}
We answer the optimal-constant question for all fractional covers satisfying a
necessary nondegeneracy condition, not only for graph incidence patterns. The
answer then feeds into a new dependent-family counting estimate.

The OpenAI formalization record for result family 161 describes a Lean theorem
asserting a finite simple host violating Sidorenko's inequality for the fixed
35-vertex, 66-edge pattern. It explicitly excludes the separate forcing
consequence from that comparator's scope \cite{OpenAILean161}. We do not use the
existence of this comparator as a proof of any new statement below.

\subsection{Relation to prior inequalities}

Fractional subadditivity and its entropy analogues are classical. Madiman and
Tetali give a broad framework of fractional inequalities and gap duality
\cite{MadimanTetali}. Ellis, Friedgut, Kindler, and Yehudayoff develop geometric
stability via information theory \cite{EllisEtAl}. Jakhar, Kurri, Chillara, and
Prabhakaran establish approximate modularity estimates controlled by a
pair-separation parameter, together with exact equality criteria
\cite{Jakhar2025}; their later work develops convex-functional extensions
\cite{Jakhar2026}.

Thus neither the principle ``a small cover gap means weak dependence'' nor the
zero-gap independence criterion is presented here as new. The contribution
studied here is the \emph{exact best constant for total defect}, expressed by
an optimization over all supports $F$, together with the incidence-span equality
classification and the exterior-algebra application. In particular, the
comparison with \cite{Jakhar2025} is quantitative and concerns a different,
global defect; Section~\ref{sec:comparison} spells out the distinction.

The quantity $m_2(G)$ in the graph specialization is the familiar 2-density
from random Ramsey theory; see, for example, the definition and discussion in
\cite{RamseyTilings}. Its identification as the constant in the span problem
is an algebraic statement. It does not by itself improve a random Ramsey
threshold or a numerical Ramsey number.

\subsection{Main conclusions and their logical separation}

The central results are Theorem~\ref{thm:cover} (all fractional covers),
Theorem~\ref{thm:span} (uniform-hypergraph spans),
Theorem~\ref{thm:equality} (strictly balanced equality cases), and
Theorem~\ref{thm:count} (dependent Lagrangian intersection counts).
The explicit source-graph improvement is Corollary~\ref{cor:residual}.
The first three results are independent of every claimed breakthrough in the
reviewed repositories. The counting theorem is also proved independently. Its
application to the particular weighted singular tail imports only the
specified residual-profile gain bound, not the entire Sidorenko conclusion.

\section{The exact fractional-cover stability constant}
\label{sec:cover}

\subsection{Definitions and the optimal-constant theorem}

Let $E$ be a finite set with $m=|E|\ge2$. Let $S_j\subseteq E$,
$j\in J$, and let $\gamma_j\ge0$ be real weights. Repetitions of the sets
$S_j$ are allowed. Empty sets and zero weights can be discarded.
The family is a \emph{fractional cover} if
\begin{equation}
 \sum_{j:e\in S_j}\gamma_j\ge1\qquad(e\in E),
 \label{eq:cover}
\end{equation}
and a \emph{fractional partition} if equality holds for every $e$.
A function $f:2^E\to\RR$ is normalized if $f(\varnothing)=0$, monotone if
$f(A)\le f(B)$ whenever $A\subseteq B$, and submodular if
\[
 f(A)+f(B)\ge f(A\cup B)+f(A\cap B).
\]
Define its total defect and cover gap by
\begin{equation}
 \Red(f)=\sum_{e\in E}f(\{e\})-f(E),\qquad
 \Gap_\gamma(f)=\sum_j\gamma_j f(S_j)-f(E).
 \label{eq:defects}
\end{equation}
Submodularity and normalization imply $\Red(f)\ge0$.
For nonempty $F\subseteq E$, put
\begin{equation}
 c_\gamma(F)=\sum_{j:S_j\cap F\ne\varnothing}\gamma_j-1.
 \label{eq:coverageexcess}
\end{equation}
Condition~\eqref{eq:cover} gives $c_\gamma(F)\ge0$.

\begin{theorem}[Exact fractional-cover constant]
\label{thm:cover}
Assume that $c_\gamma(F)>0$ whenever $|F|\ge2$, and set
\begin{equation}
 \boxed{\displaystyle
 \Kopt_\gamma=\max_{F\subseteq E,\ |F|\ge2}
 \frac{|F|-1}{c_\gamma(F)}.}
 \label{eq:K}
\end{equation}
Then every normalized monotone submodular function satisfies
\begin{equation}
 \boxed{\displaystyle
 \Red(f)\le \Kopt_\gamma\Gap_\gamma(f).}
 \label{eq:maincover}
\end{equation}
The constant is best possible, even among rank functions of one-dimensional
subspace arrangements and among entropy functions of finite random variables.
For a fractional partition, monotonicity of $f$ is unnecessary.

If $c_\gamma(F)=0$ for some $|F|\ge2$, no finite constant in
\eqref{eq:maincover} works for all normalized monotone submodular functions.
\end{theorem}

For a fractional partition, the positive-denominator condition is equivalent
to requiring that no two distinct elements have identical incidence columns
on the positive-weight sets. Indeed,
\[
 c_\gamma(\{e,f\})=1-\sum_{j:e,f\in S_j}\gamma_j.
\]
If two columns differ, their common weight is strictly less than one. Conversely,
common weight one forces their positive-weight columns to agree. Positivity on
pairs implies positivity on all larger sets.

\subsection{Unit contraction of an integral polymatroid}

The following elementary lemma is the engine of the proof. It is stated and
proved directly, so no representability assumption or matroid-expansion theorem
is needed.

\begin{lemma}[Unit contraction]
\label{lem:contraction}
Let $f$ be normalized, monotone, integer-valued, and submodular. If
$f(\{e_0\})\ge1$, define
\begin{equation}
 f'(S)=\min\{f(S),f(S\cup\{e_0\})-1\}.
 \label{eq:unitcontraction}
\end{equation}
Then $f'$ has the same four properties, $f'(E)=f(E)-1$, and
\begin{equation}
 f(S)-f'(S)=\1_{\{f(S\cup\{e_0\})=f(S)\}}.
 \label{eq:closureindicator}
\end{equation}
\end{lemma}

\begin{proof}
Write
\[
 d(S)=\1_{\{f(S\cup\{e_0\})=f(S)\}}.
\]
A nonzero marginal is a positive integer, so \eqref{eq:closureindicator} holds.
Diminishing returns shows that $d(A)\le d(B)$ for $A\subseteq B$.
Also $d(\varnothing)=0$, because $f(\{e_0\})\ge1$, and $d(E)=1$.

For monotonicity, the only nontrivial case is $A\subseteq B$ with
$d(A)=0$ and $d(B)=1$. If $f(A)=f(B)$, monotonicity and $d(B)=1$ would give
\[
 f(A\cup\{e_0\})\le f(B\cup\{e_0\})=f(B)=f(A),
\]
contradicting $d(A)=0$. Consequently $f(B)-f(A)\ge1$, and
$f'(B)-f'(A)=f(B)-f(A)-1\ge0$.

For submodularity set $I=A\cap B$ and $U=A\cup B$. We need
\[
 f(A)+f(B)-f(U)-f(I)\ge d(A)+d(B)-d(U)-d(I).
\]
The right side is positive only if $d(A)=d(B)=d(U)=1$ and $d(I)=0$;
otherwise it is at most zero. In that exceptional case, applying submodularity
to $A\cup\{e_0\}$ and $B\cup\{e_0\}$ gives
\[
 f(A)+f(B)\ge f(U)+f(I\cup\{e_0\})\ge f(U)+f(I)+1.
\]
This proves the required inequality. Normalization, integrality, and
$f'(E)=f(E)-1$ have already been checked. Monotonicity and normalization also
give nonnegativity.
\end{proof}

\subsection{Proof of the exact constant}

\begin{proof}[Proof of Theorem~\ref{thm:cover}]
First suppose that $f$ is integer-valued and monotone. We induct on $f(E)$.
If $f(E)=0$, monotonicity forces $f$ to vanish. Otherwise subadditivity ensures
that some $e_0$ has $f(\{e_0\})>0$. Apply Lemma~\ref{lem:contraction} and put
\[
 F=\{e:f(\{e,e_0\})=f(\{e\})\},\qquad
 U=\{j:f(S_j\cup\{e_0\})=f(S_j)\}.
\]
The set $F$ contains $e_0$. If $S_j$ meets $F$, diminishing returns makes the
marginal of $e_0$ over $S_j$ zero, so $j\in U$. Therefore
\begin{align}
 \Red(f)-\Red(f')&=|F|-1,\label{eq:Rincrement}\\
 \Gap_\gamma(f)-\Gap_\gamma(f')&=\sum_{j\in U}\gamma_j-1
 \ge c_\gamma(F).\label{eq:Gincrement}
\end{align}
If $|F|=1$, the first increment is zero and the second is nonnegative.
If $|F|\ge2$, definition~\eqref{eq:K} yields
\[
 |F|-1\le \Kopt_\gamma c_\gamma(F)
 \le \Kopt_\gamma\bigl(\Gap_\gamma(f)-\Gap_\gamma(f')\bigr).
\]
Adding the induction hypothesis for $f'$ proves the result.

Scaling proves the inequality for rational-valued monotone submodular
functions. For completeness, real-valued functions can be approximated without
leaving the monotone submodular cone. Put
\[
 h(S)=(2m+1)|S|-|S|^2.
\]
Its nontrivial monotonicity increments are at least two, and its submodularity
slack on incomparable $A,B$ is $2|A\setminus B||B\setminus A|\ge2$.
Thus $f+\varepsilon h$ can be approximated coordinatewise by a normalized
rational-valued function within $\varepsilon/4$, while preserving all relevant
inequalities. Comparable-set submodularity inequalities are identities.
Letting $\varepsilon\downarrow0$ proves the real case by continuity.

When the weights form a fractional partition, adding a modular function
$M|S|$ changes neither defect in \eqref{eq:defects}. For any normalized
submodular $f$, a sufficiently large $M$ makes $f(S)+M|S|$ monotone. The already
proved case therefore removes the monotonicity hypothesis.

To prove optimality, fix $F\subseteq E$, $|F|\ge2$, and use the function
\begin{equation}
 f_F(S)=\1_{\{S\cap F\ne\varnothing\}}.
 \label{eq:rankone}
\end{equation}
It is normalized, monotone, and submodular, and
\[
 \Red(f_F)=|F|-1,\qquad \Gap_\gamma(f_F)=c_\gamma(F).
\]
It is realized by assigning one common nonzero line to the elements of $F$
and the zero subspace to the others. It is also, up to multiplication by the
entropy of a nonconstant finite random variable $Z$, the entropy function
obtained by setting $X_e=Z$ for $e\in F$ and making the remaining variables
constant. Taking a maximizing $F$ proves sharpness. A zero denominator gives a
positive total defect and a zero cover gap, excluding every finite constant.
\end{proof}

\begin{remark}[What the contraction proof does not assert]
Repeated contractions do not express an arbitrary submodular function as a sum
of coverage functions. They decompose the \emph{two particular defects} into
increments with controlled ratios. In particular, the proof neither assumes
that every polymatroid is linearly representable nor makes such a conclusion.
\end{remark}

\subsection{A finite network for the optimal constant}

The support formula is explicit but contains exponentially many subsets. Its
comparison with a prescribed rational number has a useful network formulation.

\begin{proposition}[A cut formulation]
\label{prop:cut}
For an explicitly given rational fractional cover and a rational $\kappa>0$,
the assertion $\Kopt_\gamma>\kappa$ is decidable using at most
$\binom{m}{2}$ minimum-cut computations on networks with $m+|J|+2$ vertices.
\end{proposition}

\begin{proof}
The assertion is equivalent to the existence of $F$, $|F|\ge2$, with
\[
 |F|-\kappa\sum_{j:S_j\cap F\ne\varnothing}\gamma_j>1-\kappa.
\]
Fix a pair of elements that must belong to $F$. Create a source $s$, a sink
$t$, one vertex for each $e\in E$, and one vertex for each $j\in J$.
Use capacity one on $s\to e$, capacity $\kappa\gamma_j$ on $j\to t$, and a
capacity $M$ on $e\to j$ whenever $e\in S_j$, where $M$ exceeds the sum of all
finite capacities. Add capacity-$M$ source arcs to the two forced elements.
A minimum cut cannot cut any $M$-arc. Its source-side element vertices form a
set $F$ containing the forced pair, and every row meeting $F$ must also lie on
the source side. Extra row vertices cannot lower the cut. Its finite capacity
is consequently
\[
 m-|F|+\kappa\sum_{j:S_j\cap F\ne\varnothing}\gamma_j.
\]
The required strict inequality is equivalent to a cut value smaller than
$m-1+\kappa$. Enumerating the forced pair covers all $|F|\ge2$.
\end{proof}

This is an algorithmic consequence of the exact constant, not a claim of a new
minimum-cut algorithm. The supplied verifier uses direct exact enumeration for
the small certificates, avoiding a dependency on a network-flow library.

\section{Entropy, determinants, and the lower cover gap}
\label{sec:entropy}

\subsection{Optimal control of total correlation}

Let $(X_e)_{e\in E}$ be finite-valued random variables and write $H(X_S)$ for
joint entropy, using natural logarithms. Then
\begin{equation}
 \Red(H)=\sum_e H(X_e)-H(X_E)
 =\KL\!\left(P_{X_E}\,\middle\|\,\bigotimes_e P_{X_e}\right).
 \label{eq:TC}
\end{equation}
The entropy set function is normalized, monotone, and submodular, so
Theorem~\ref{thm:cover} gives the following sharp statement.

\begin{corollary}[Sharp entropy stability]
\label{cor:entropy}
For every fractional cover satisfying the positive-denominator condition,
\begin{equation}
 \KL\!\left(P_{X_E}\,\middle\|\,\bigotimes_e P_{X_e}\right)
 \le\Kopt_\gamma\left(\sum_j\gamma_j H(X_{S_j})-H(X_E)\right).
 \label{eq:entropybound}
\end{equation}
The coefficient cannot be decreased, even for binary-valued variables.
In particular, with $g$ denoting the expression in parentheses,
\begin{equation}
 \TV\!\left(P_{X_E},\bigotimes_eP_{X_e}\right)
 \le\sqrt{\Kopt_\gamma g/2}.
 \label{eq:TV}
\end{equation}
\end{corollary}

\begin{proof}
Only the final estimate needs an additional argument. For probability measures
$P,Q$, let $A=\{x:P(x)\ge Q(x)\}$. The log-sum inequality gives
$\KL(P\|Q)\ge d(P(A)\|Q(A))$, where $d$ is binary relative entropy.
As a function of its first argument, $d(p\|q)$ has second derivative
$1/(p(1-p))\ge4$ and value and derivative zero at $p=q$.
Hence $d(p\|q)\ge2(p-q)^2$. Since $P(A)-Q(A)=\TV(P,Q)$, this proves
\eqref{eq:TV}. Sharpness of the coefficient in \eqref{eq:entropybound} follows
from the common binary variable construction in \eqref{eq:rankone}; no claim
of sharpness of the square-root bound for every fixed cover is needed.
\end{proof}

The same entropy inequality holds conditionally on another random variable:
apply it to each conditional distribution and average. Its constant is still
sharp because the conditioning variable may be constant. This extension does
not require a linear or Gaussian model of the dependence.

\subsection{A sharp determinant inequality}

\begin{corollary}[Positive-definite matrices]
\label{cor:determinants}
Let $(S_j,\gamma_j)$ be a fractional partition with finite $\Kopt_\gamma$.
For every real positive-definite matrix $\Sigma$ indexed by $E$,
\begin{equation}
 \boxed{\displaystyle
 \frac{\prod_{e\in E}\Sigma_{ee}}{\det\Sigma}
 \le
 \left(\frac{\prod_j(\det\Sigma_{S_j})^{\gamma_j}}
 {\det\Sigma}\right)^{\Kopt_\gamma}.}
 \label{eq:detbound}
\end{equation}
Here the determinant of the empty principal submatrix is one. The exponent is
optimal even if one restricts to covariance matrices of nonsingular Gaussian
vectors.
\end{corollary}

\begin{proof}
The set function $f(S)=\log\det\Sigma_S$, with $f(\varnothing)=0$, is
submodular. One direct proof conditions on the common block $C=A\cap B$ by
Schur complementation. The remaining positive-definite block matrix on
$A\setminus B$ and $B\setminus A$ has determinant at most the product of its
two diagonal-block determinants: take a second Schur complement and use
$0<M_{QQ}-M_{QP}M_{PP}^{-1}M_{PQ}\le M_{QQ}$.
Taking determinants gives the required inequality for $f$.
This $f$ need not be monotone, which is why the fractional-partition version of
Theorem~\ref{thm:cover} is used. Exponentiating its conclusion yields
\eqref{eq:detbound}.

For sharpness choose a maximizing support $F$, put
$\Sigma_F=\mathbf 1\mathbf 1^{\mathsf T}+\varepsilon I$, and use an independent
identity block outside $F$. For every $S$ with $r=|S\cap F|\ge1$,
\[
 \det\Sigma_S=\varepsilon^{r-1}(r+\varepsilon).
\]
Consequently, as $\varepsilon\downarrow0$,
\begin{align*}
 \log\frac{\prod_e\Sigma_{ee}}{\det\Sigma}
 &=(|F|-1)\log(1/\varepsilon)+O(1),\\
 \log\frac{\prod_j(\det\Sigma_{S_j})^{\gamma_j}}{\det\Sigma}
 &=c_\gamma(F)\log(1/\varepsilon)+O(1).
\end{align*}
In the second identity we used
$\sum_j\gamma_j|S_j\cap F|=|F|$. Their ratio tends to
$\Kopt_\gamma$. Every matrix in the approximation is positive definite.
\end{proof}

\subsection{The exact lower-gap constant}

For a fractional partition define
\[
 \Gap^-_\gamma(f)=f(E)-\sum_j\gamma_j
 \bigl(f(E)-f(E\setminus S_j)\bigr).
\]
This is the lower fractional-subadditivity gap in conditional notation.
The standard gap duality \cite{MadimanTetali} has an exact-constant consequence.

\begin{corollary}[Optimal comparison with the lower gap]
\label{cor:lower}
For a separating fractional partition, put
\begin{equation}
 \Kopt^-_\gamma=\max_{|F|\ge2}
 \frac{|F|-1}{1-\sum_{j:F\subseteq S_j}\gamma_j}.
 \label{eq:Kminus}
\end{equation}
Then every normalized submodular function satisfies
$\Red(f)\le\Kopt^-_\gamma\Gap^-_\gamma(f)$, and this coefficient is sharp.
\end{corollary}

\begin{proof}
Let $w=\sum_j\gamma_j$. Separation implies $w>1$. The complementary sets
$T_j=E\setminus S_j$, with weights $\bar\gamma_j=\gamma_j/(w-1)$, form a
fractional partition, and direct expansion gives
\[
 \Gap^-_\gamma(f)=(w-1)\Gap_{\bar\gamma}(f).
\]
For $|F|\ge2$,
\[
 (w-1)c_{\bar\gamma}(F)
 =1-\sum_{j:F\subseteq S_j}\gamma_j>0.
\]
Theorem~\ref{thm:cover} proves the inequality. The rank-one functions $f_F$
again give equality for a maximizing support.
\end{proof}

For the star partition of a simple graph with $m\ge2$ edges and maximum degree
$d_{\max}$, normalized by weights $1/2$, formula~\eqref{eq:Kminus} simplifies to
\begin{equation}
 \Kopt^-_\gamma=\max\{m-1,\,2(d_{\max}-1)\}.
 \label{eq:graphlower}
\end{equation}
Indeed, a collection of at least two distinct edges is contained in at most one
star. A support contained in a star has denominator $1/2$; every other support
has denominator one. If the whole graph is a star, the star term dominates;
otherwise the whole edge set attains $m-1$. This also shows why the optimal
upper-gap and lower-gap constants need not resemble one another.

\subsection{Comparison with pair-separation stability}
\label{sec:comparison}

For a separating fractional partition, \cite[Theorem 3]{Jakhar2025} controls
$f(\{e\})+f(E\setminus\{e\})-f(E)$ using the gap and
\[
 \sigma=\min_{e\ne f}\sum_{j:e\in S_j,\ f\notin S_j}\gamma_j.
\]
A chain-rule summation of those pointwise estimates gives the valid comparison
\begin{equation}
 \Red(f)\le\frac{m-1}{\sigma}\Gap_\gamma(f).
 \label{eq:paircomparison}
\end{equation}
For example, each term
$f(\{e\})-f(\{e\}\mid\text{earlier elements})$ is bounded above by the
corresponding single-element versus all-other-elements defect; the first term
in the chain is zero. This derives \eqref{eq:paircomparison} without attributing
to the cited paper a statement in precisely this form.

For a graph with adjacent edges, the star partition has $\sigma=1/2$.
The new exact comparison will be
$\Red(f)\le 2m_2(G)\Gap_\gamma(f)$, replacing $2(m-1)$ by $2m_2(G)$.
For the 33-edge source graph, the improvement in this particular derived global
bound is a factor of eleven. This does not imply that every conclusion of
\cite{Jakhar2025} is weakened by that factor; its pointwise estimates address
additional quantities.

\section{Incidence spans, directness repair, and equality}
\label{sec:spans}

\subsection{Uniform hypergraphs and the sharp span inequality}

Let $\mathcal H=(V,E)$ be a finite simple $a$-uniform hypergraph with at least
two edges. The word simple means that different edges are different $a$-element
subsets. Isolated vertices may be removed. For $F\subseteq E$ let
$V(F)=\bigcup_{e\in F}e$, and define
\begin{equation}
 \mu_a(\mathcal H)=\max_{F\subseteq E,\ |F|\ge2}
 \frac{|F|-1}{|V(F)|-a}.
 \label{eq:adensity}
\end{equation}
All denominators are positive. In a graph, this is the usual $m_2(G)$.

Assign finite-dimensional subspaces $H_e$ of a common vector space over an
arbitrary field. Put
\begin{equation}
 \begin{gathered}
 d_e=\dim H_e,\qquad t=\sum_e d_e,\qquad
 K=\sum_eH_e,\quad k=\dim K,\\
 K_v=\sum_{e\ni v}H_e,\qquad k_v=\dim K_v,\qquad
 R=t-k,\qquad \Delta=\sum_v k_v-a k.
 \end{gathered}
 \label{eq:spanparameters}
\end{equation}
The quantity $R$ is the dimension of the kernel of the addition map
$\bigoplus_eH_e\to K$. It measures all linear relations, not merely the number
of pairs with nonzero intersection.

\begin{theorem}[Sharp incidence-span inequality]
\label{thm:span}
For every such arrangement,
\begin{equation}
 \boxed{\displaystyle R\le\mu_a(\mathcal H)\Delta.}
 \label{eq:spanmain}
\end{equation}
The coefficient is optimal for each fixed hypergraph. In particular,
$\Delta=0$ if and only if the $H_e$ form a direct sum. For graphs this reads
\begin{equation}
 \sum_v\dim K_v-2\dim K
 \ge\frac{\sum_e\dim H_e-\dim K}{m_2(G)}.
 \label{eq:graphspan}
\end{equation}
\end{theorem}

\begin{proof}
For $F\subseteq E$, the function $f(F)=\dim\sum_{e\in F}H_e$ is normalized,
monotone, and submodular. The stars $S_v=\{e:v\in e\}$, with weights $1/a$,
form a fractional partition. Here
\[
 \Gap_\gamma(f)=\Delta/a,\qquad
 c_\gamma(F)=(|V(F)|-a)/a,\qquad
 \Kopt_\gamma=a\mu_a(\mathcal H).
\]
Theorem~\ref{thm:cover} proves the inequality and sharpness.
If $\Delta=0$, then $R=0$, which is directness. Conversely, a direct family has
$k=\sum_e d_e$ and $k_v=\sum_{e\ni v}d_e$, so $\Delta=0$.
\end{proof}

\subsection{An exact flag identity over every field}

A direct linear proof supplies additional information needed for equality.
Choose a nonzero line $L$ lying in at least one $H_e$, and quotient $K$ by $L$.
Use primes for the images and their dimensions. Set
\[
 F_L=\{e:L\subseteq H_e\},\qquad U_L=\{v:L\subseteq K_v\}.
\]
Then $F_L\ne\varnothing$, $V(F_L)\subseteq U_L$, and rank-nullity gives
\begin{equation}
 R-R'=|F_L|-1,\qquad
 \Delta-\Delta'=|U_L|-a.
 \label{eq:linearstep}
\end{equation}
Repeating until the ambient span becomes zero yields a flag with $k$ steps and
\begin{equation}
 \boxed{\displaystyle
 R=\sum_{\ell=1}^k(|F_\ell|-1),\qquad
 \Delta=\sum_{\ell=1}^k(|U_\ell|-a).}
 \label{eq:flag}
\end{equation}
The supports and visible vertex sets are taken in the \emph{successive
quotients}, not all in the original arrangement. At each step,
\[
 |F_\ell|-1\le\mu_a(|V(F_\ell)|-a)
 \le\mu_a(|U_\ell|-a)
\]
when $|F_\ell|\ge2$; the singleton case contributes zero on the left and a
nonnegative quantity on the right. This proves \eqref{eq:spanmain} again.

\begin{corollary}[Equality certificate by a flag]
\label{cor:flag}
Equality in \eqref{eq:spanmain} holds if and only if every step of one such flag
has one of the following forms: a singleton support with $|U_\ell|=a$, or a
support maximizing \eqref{eq:adensity} with $U_\ell=V(F_\ell)$.
If equality holds, this property holds for every admissible flag.
\end{corollary}

\begin{proof}
The total slack is the sum of the nonnegative step slacks in
\eqref{eq:linearstep}. A step slack is zero precisely under the stated
conditions. This reasoning does not depend on which admissible flag was chosen.
\end{proof}

\subsection{The exact cost of repairing to a direct sum}

\begin{theorem}[Optimal directness repair]
\label{thm:repair}
Among all families $H'_e\subseteq H_e$ forming a direct sum, the minimum of
$\sum_e\codim_{H_e}H'_e$ is exactly $R$. In particular a direct family can be
obtained by deleting at most $\mu_a(\mathcal H)\Delta$ dimensions in total,
and the resulting total dimension can be $k$.
\end{theorem}

\begin{proof}
A direct family of subspaces of $K$ has total dimension at most $k$, so its
number of deleted dimensions is at least $t-k=R$. Conversely, choose a basis
of each $H_e$, scan the union of these bases, and retain a vector precisely
when it enlarges the span of the previously retained vectors. The retained
vectors form a basis of $K$. Assign each retained vector to its original edge
and let $H'_e$ be the span of the vectors assigned to $e$. Their total union
is a basis, so the family is direct and has total dimension $k$.
The final bound follows from Theorem~\ref{thm:span}.
\end{proof}

This is a dimension-edit statement, not a claim that a near-extremal family has
a common core after a similarly small edit. The latter, different stability
problem is discussed in Section~\ref{sec:questions}.

\subsection{Equality for strictly balanced hypergraphs}

\begin{definition}
A simple $a$-uniform hypergraph with $m\ge2$ edges is \emph{strictly
$a$-balanced} if its full edge set is the unique maximizer in
\eqref{eq:adensity}. Thus
$\mu_a=(m-1)/(|V|-a)$ and every proper support of at least two edges has
strictly smaller density.
\end{definition}

\begin{theorem}[Common-core classification]
\label{thm:equality}
Suppose $\mathcal H$ is strictly $a$-balanced. Equality
$R=\mu_a(\mathcal H)\Delta$ holds if and only if there are subspaces $C,P_e$
such that
\begin{equation}
 H_e=C\oplus P_e\quad(e\in E),\qquad
 K=C\oplus\bigoplus_{e\in E}P_e.
 \label{eq:commoncore}
\end{equation}
Necessarily $C=\bigcap_eH_e$. In particular, if the total intersection is zero,
equality is equivalent to directness.
\end{theorem}

\begin{proof}
Suppose equality holds. Take any nonzero vector in $H_e\cap H_f$ for $e\ne f$
and begin a flag by its line. Corollary~\ref{cor:flag} and strict balance force
$F_L=E$. Hence every pairwise intersection equals
$C=\bigcap_e H_e$.

Fix an edge $e_0$ and quotient by $H_{e_0}$, using successive one-dimensional
quotients inside its current image. Equality is preserved at every step,
because both the step slack and the remaining slack are nonnegative and their
sum is zero. In the final quotient the $e_0$-space is zero. No nonzero line
there can have support $E$. Thus a further flag can contain only singleton
supports, so its redundancy is zero. Writing $c=\dim C$ and comparing total
edge dimensions before and after the quotient gives
\[
 0=R_{\mathrm{quotient}}
 =R-\sum_{e\ne e_0}\dim(H_e\cap H_{e_0})
 =R-(m-1)c.
\]
Choose any complement $P_e$ to $C$ in $H_e$. The addition map from
$C\oplus\bigoplus_eP_e$ onto $K$ has domain dimension
$c+\sum_e(d_e-c)=t-(m-1)c=k$, so it is an isomorphism. This proves
\eqref{eq:commoncore}.

Conversely, an arrangement of the form \eqref{eq:commoncore} satisfies
\[
 R=(m-1)c,\qquad \Delta=(|V|-a)c.
\]
Their ratio is $\mu_a$, proving equality.
\end{proof}

\section{Ramsey density and the source complex}
\label{sec:density}

\subsection{Several exact graph constants}

For graphs with at least two edges, isolated vertices do not affect $m_2$.
The following examples make the optimal constant concrete.
\begin{center}
\begin{tabular}{@{}lll@{}}
\toprule
Graph & $m_2(G)$ & Equality structure covered above\\
\midrule
$K_n$, $n\ge3$ & $(n+1)/2$ & Common core and private summands\\
$C_\ell$, $\ell\ge3$ & $(\ell-1)/(\ell-2)$ & Common core and private summands\\
Forest with adjacent edges & $1$ & May have several maximizing supports\\
Matching with at least two edges & $1/2$ & May have several maximizing supports\\
Maximal planar graph, $n\ge4$ & $3-1/(n-2)$ & Common core and private summands\\
\bottomrule
\end{tabular}
\end{center}

For $K_n$, any support on $s$ vertices has at most $\binom{s}{2}$ edges, so
its density is at most $(s+1)/2$; equality at the maximum requires the entire
$K_n$. A proper support of a cycle is a forest and has density at most one,
whereas the complete cycle has density $(\ell-1)/(\ell-2)>1$.
A forest has at most $v-1$ edges on $v$ vertices, and an adjacent pair attains
density one. A matching support of $r\ge2$ edges has $2r$ vertices and density
$(r-1)/(2r-2)=1/2$.

For a simple planar graph on $s\ge3$ vertices, the familiar edge bound
$e\le3s-6$ follows by Euler's formula and face-length counting, with disconnected
cases obtained by adding planar connecting edges if necessary. A support in a
maximal planar graph on $n$ vertices therefore has density at most
\[
 3-\frac1{s-2}\le3-\frac1{n-2}.
\]
The second inequality is strict for $s<n$, and a proper support on all $n$
vertices loses at least one edge. Thus maximal planar graphs are strictly
2-balanced, giving the last line of the table.

These calculations explain the appearance of the same density parameter as in
random Ramsey theory without relying on a probabilistic threshold theorem.
The shared parameter is useful, but an implication between the two theories
would require additional arguments.

\subsection{The 13-point graph}

The Sidorenko manuscript specifies 22 triples on 13 points. Its bipartite
pattern has one part consisting of points and the other of triples; incidence
creates 35 vertices and 66 edges. The graph relevant to the span calculation
is a \emph{different} graph: the 13-vertex point graph $G_*$, whose edges are
pairs contained in a triple. It has 33 edges. Confusing these two graphs would
give the wrong density and the wrong constant.

Here is the triple system from the source \cite{OpenAISidorenko}:
\begin{equation}
\begin{array}{rrrr}
013&019&023&029\\
1\,3\,10&1\,7\,10&1\,7\,12&1\,9\,12\\
234&249&3\,4\,11&3\,10\,11\\
489&4\,8\,11&567&5\,6\,11\\
5\,7\,10&5\,10\,11&6\,7\,12&6\,8\,11\\
6\,8\,12&8\,9\,12&&
\end{array}
\label{eq:faces}
\end{equation}
Entries such as $013$ denote the set $\{0,1,3\}$, not a decimal integer.
The degree sequence is
\begin{equation}
 (4,6,4,6,5,4,5,5,5,6,5,6,5).
 \label{eq:degrees}
\end{equation}
Appendix~\ref{app:rotation} gives a rotation system whose face cycles are exactly
these 22 triangles. It is connected and has Euler characteristic
$13-33+22=2$, so it is a spherical embedding. Thus $G_*$ is maximal planar.
Alternatively, the included exact subset certificate directly checks its
density on every vertex subset of size at least three.

\begin{corollary}[Optimal source-graph span bound]
\label{cor:sourcegraph}
For subspaces indexed by the 33 edges of $G_*$,
\begin{equation}
 m_2(G_*)=\frac{32}{11},\qquad
 \boxed{32\Delta\ge11R.}
 \label{eq:3211}
\end{equation}
The coefficient is optimal. Equality holds precisely for the common-core
arrangements of Theorem~\ref{thm:equality}.
\end{corollary}

\begin{proof}
Since $G_*$ is maximal planar with 13 vertices,
$m_2(G_*)=3-1/11=32/11$, and it is strictly 2-balanced. Apply
Theorems~\ref{thm:span} and~\ref{thm:equality}. Assigning the same line to every
edge gives $R=32$ and $\Delta=11$, establishing sharpness concretely.
\end{proof}

For integer dimensions a useful form is
\begin{equation}
 \Delta\ge\left\lceil\frac{11R}{32}\right\rceil.
 \label{eq:ceilbound}
\end{equation}
The original extra-unit assertion detects only whether $R$ is zero. In
contrast, \eqref{eq:ceilbound} forces increasingly large excess as the number
of relations grows, which is exactly what is needed to offset the entropy of
choosing a relation subspace in the next section.

\section{Counting dependent Lagrangian intersection families}
\label{sec:counting}

\subsection{The event to be counted}

Let $G=(V,E)$ be any fixed finite simple graph. Prescribe nonnegative integers
$h_e$. Let $L_v$, $v\in V$, be independent uniformly distributed Lagrangian
subspaces of a $2D$-dimensional nondegenerate symplectic space over $\Fq$.
A Lagrangian has dimension $D$ and is totally isotropic.
We count families
\begin{equation}
 H_e\subseteq L_u\cap L_v\quad(e=\{u,v\}),\qquad \dim H_e=h_e.
 \label{eq:containment}
\end{equation}
Containment, rather than equality, is intentional: it permits an upper bound
for actual intersection profiles while avoiding unnecessary exactness
constraints.

Write
\begin{equation}
 t=\sum_e h_e,\quad t_v=\sum_{e\ni v}h_e,\quad
 k=\dim\sum_eH_e,\quad k_v=\dim\sum_{e\ni v}H_e,
 \quad R=t-k,\quad \Delta=\sum_v k_v-2k.
 \label{eq:countparameters}
\end{equation}
Define $b(x)=x(x+1)/2$. All constants in this section may depend on the fixed
graph and prescribed dimensions but can be chosen uniformly in $q$ and $D$.
Empty or dimensionally impossible strata contribute zero.

\subsection{Two elementary symplectic counts}

\begin{lemma}[Lagrangian containment probability]
\label{lem:lagcount}
If $A$ is a fixed isotropic $s$-space, $s\le D$, and $L$ is a uniform
Lagrangian, then
\begin{equation}
 \Pp(A\subseteq L)=\prod_{i=D-s+1}^{D}(q^i+1)^{-1}
 \le q^{-Ds+\binom{s}{2}}.
 \label{eq:lagprob}
\end{equation}
\end{lemma}

\begin{proof}
Let $N_D$ be the number of Lagrangians in dimension $2D$.
Every line is isotropic. Count pairs consisting of a line and a Lagrangian
containing it. The quotient of the line's orthogonal complement by the line is
symplectic of dimension $2(D-1)$, so
\[
 N_D\frac{q^D-1}{q-1}=\frac{q^{2D}-1}{q-1}N_{D-1}.
\]
Thus $N_D=(q^D+1)N_{D-1}$ and $N_D=\prod_{i=1}^D(q^i+1)$.
The quotient $A^\perp/A$ similarly identifies the Lagrangians containing $A$
with those in dimension $2(D-s)$, proving the product formula.
The exponent in the bound is minus
$\sum_{i=D-s+1}^D i=-Ds+\binom{s}{2}$.
\end{proof}

\begin{lemma}[Embedding count with linear Gram constraints]
\label{lem:gram}
Let $P$ be a $k$-dimensional vector space and let
$\mathcal L\subseteq\Lambda^2P$ have dimension $\ell$.
The number of injective linear maps $\phi:P\to\Fq^{2D}$ whose pulled-back
symplectic form vanishes on $\mathcal L$ is at most
\begin{equation}
 q^{2Dk-\ell}.
 \label{eq:gramcount}
\end{equation}
\end{lemma}

\begin{proof}
There are $q^{\binom{k}{2}-\ell}$ alternating forms on $P$ annihilating
$\mathcal L$. Fix one such form and a basis of $P$. Having chosen $i-1$
independent image vectors, prescribing the pairings of the next vector with
them imposes $i-1$ independent linear equations, by nondegeneracy of the
ambient form. There are at most $q^{2D-(i-1)}$ choices; excluding dependent
choices can only decrease this number. Multiplying and then summing over the
forms gives \eqref{eq:gramcount}. If $k>2D$, no injection exists.
The proof uses alternating forms and remains valid in characteristic two.
\end{proof}

\subsection{Relations and exterior constraints}

Set
\[
 U=\bigoplus_{e\in E}U_e,\quad U_e=\Fq^{h_e},\qquad
 U_v=\bigoplus_{e\ni v}U_e.
\]
For a family of based subspaces $H_e$, the basis vectors define a map
$\Phi:U\to\Fq^{2D}$. Its kernel $N$ has dimension $R$, and its quotient
$P=U/N$ has dimension $k$. Let $\pi:U\to P$ be the quotient map. Put
\begin{equation}
 \mathcal L_0=\sum_v\Lambda^2U_v\subseteq\Lambda^2U,
 \qquad
 \mathcal L=(\Lambda^2\pi)(\mathcal L_0).
 \label{eq:exteriorconstraints}
\end{equation}
Local isotropy of all $\Phi(U_v)$ is exactly the condition that the pulled-back
alternating form vanish on $\mathcal L$.
Because the graph is simple, distinct edges share at most one endpoint.
Counting basis wedges therefore gives
\begin{equation}
 \ell_0:=\dim\mathcal L_0
 =\sum_v\binom{t_v}{2}-\sum_e\binom{h_e}{2}.
 \label{eq:ellzero}
\end{equation}
A wedge inside a single edge block is counted twice in the first sum; a wedge
between adjacent distinct edge blocks is counted once, and a wedge between
disjoint edges is not counted. This proves the identity without an independence
assumption on the eventual image subspaces.

Let $\ell=\dim\mathcal L$ and define
\begin{equation}
 B_N=\ell_0-\ell,\qquad
 A_N=\sum_v\left[\binom{t_v}{2}-\binom{k_v}{2}\right].
 \label{eq:AB}
\end{equation}
The second quantity is nonnegative. Choosing a complement to $N$ in $U$ shows
\begin{equation}
 \dim\ker(\Lambda^2\pi)=Rk+\binom R2,
 \qquad 0\le B_N\le Rk+\binom R2.
 \label{eq:exteriorkernel}
\end{equation}
Indeed, the kernel is the sum of $N\wedge P$ and $\Lambda^2N$ in such a
splitting. These dimensions, not an assertion of independence of arbitrary
orthogonality equations, will control the dependent case.

\subsection{The dependent-family counting theorem}

\begin{theorem}[Relation-sensitive Lagrangian count]
\label{thm:count}
For fixed $h_e,k,k_v$, the probability that a family satisfying
\eqref{eq:containment} exists with those ranks is at most
\begin{equation}
 C\,q^{-\sum_e b(h_e)-D\Delta+2Rk+\binom R2-A_N}.
 \label{eq:maincount}
\end{equation}
Here $A_N$ is determined by the prescribed $k_v$ and hence is constant on the
stratum. A more precise bound on a stratum where $B_N=B$ is fixed is
\begin{equation}
 C\,q^{-\sum_e b(h_e)-D\Delta+Rk+B-A_N}.
 \label{eq:finecount}
\end{equation}
The constant $C$ depends only on the finite graph and the prescribed edge
dimensions. In particular, it can be chosen independently of $D$ and $q$.
\end{theorem}

\begin{proof}
There are at most $C_t q^{R(t-R)}=C_tq^{Rk}$ possible relation spaces
$N\subseteq U$ of dimension $R$; this is the usual Gaussian-binomial bound,
also obtained by counting full-rank matrices and dividing by ordered bases.
For each fixed $N$, discard it if an edge block is not injective in $P=U/N$
or if it has the wrong local ranks. Lemma~\ref{lem:gram} bounds the remaining
injective maps $P\to\Fq^{2D}$ by $q^{2Dk-\ell}$.

A family of subspaces $H_e$ is represented by exactly
$\prod_e|\mathrm{GL}(h_e,q)|$ choices of its ordered edge bases. These choices
can change $N$, but preserve every rank and the dimensions of the exterior
constraint spaces. We have summed over all such $N$, so division by this
multiplicity is legitimate. For fixed $h_e$ it is bounded below by a positive
constant times $q^{\sum_eh_e^2}$.

For each resulting locally isotropic family, independence of the Lagrangians
and Lemma~\ref{lem:lagcount} bound its containment probability by
\[
 q^{-D\sum_v k_v+\sum_v\binom{k_v}{2}}.
\]
The expectation of the number of such contained families is therefore at most
a constant times $q$ to the exponent
\begin{align*}
 Rk+2Dk-\ell-\sum_eh_e^2
 -D\sum_v k_v+\sum_v\binom{k_v}{2}
 &= -D\Delta+Rk-\sum_eh_e^2-\ell
       +\sum_v\binom{k_v}{2}\\
 &= -\sum_e b(h_e)-D\Delta+Rk+B_N-A_N.
\end{align*}
The last equality uses \eqref{eq:ellzero} and
$h^2-\binom h2=b(h)$. The event that a family exists has probability at most
this expectation. This proves \eqref{eq:finecount}; using
\eqref{eq:exteriorkernel} proves \eqref{eq:maincount}.
Only boundedly many integer exterior ranks can occur, so summing over $B_N$
changes the constant, not the exponent bound.
\end{proof}

\begin{remark}[Why the ambient span need not be isotropic]
Only the local spans $K_v$ must be isotropic. Pairs of edge spaces with no common
endpoint need not be mutually orthogonal. The Gram-space proof retains exactly
the local equations through $\mathcal L$ and never imposes isotropy on the
whole $k$-dimensional ambient span. Replacing that span by a globally isotropic
space would be an unjustified restriction of the event being bounded.
\end{remark}

\begin{corollary}[Weights and direct families]
\label{cor:weighted}
If a nonnegative weight $W$ is at most $C_0q^g$ on the event in
Theorem~\ref{thm:count}, its weighted contribution is bounded by the same
expressions with $g$ added to the exponent and $C_0$ absorbed into the constant.
When $R=0$, one has $\Delta=A_N=B_N=0$, and the unweighted bound reduces to
$Cq^{-\sum_e b(h_e)}$.
\end{corollary}

The direct-family exponent agrees with the direct orthogonality cancellation
in the source \cite[\src{lem:tail-residual-counts}]{OpenAISidorenko}.
The new point is a controlled extension of that cancellation to dependent
families, with explicit payment for choosing relations and for losing exterior
constraints.

\section{An explicit improvement of the residual singular tail}
\label{sec:application}

\subsection{The precise imported interface}

The source's earlier profile analysis isolates possible residual strata with
\begin{equation}
 h_e\in\{0,2\},\qquad c_j=0,
 \label{eq:residualprofile}
\end{equation}
where $c_j$ is the dimension common to the three point Lagrangians on a face.
For a residual stratum with $m$ nonzero edge spaces, the center-integrated
nonnegative weight is bounded by $O_D(q^{2m})$. This is the gain
$\frac23\sum_e b(h_e)=2m$ stated in the source's residual-profile argument
\cite[\src{sections/tail.tex}, \src{lem:tail-residue} and
\src{lem:tail-residual-counts}]{OpenAISidorenko}.

For clarity, our new estimate has the following standalone interpretation:
\emph{for any nonnegative weight satisfying this pointwise bound on these
intersection strata, the conclusion below follows}. Neither the finite-field
rank-layer limit nor the counterexample theorem is needed to prove that
statement. Connecting the weight to the particular source construction uses
the just-specified gain bound.

\subsection{The rank-dependent exponent}

There are at most 33 nonzero edges, so $1\le m\le33$ on a singular residual
stratum. Since $t=2m$, one has $k=2m-R$. Dropping the helpful nonnegative term
$A_N$ from Theorem~\ref{thm:count} and adding the gain $2m$ gives
\begin{align}
 -3m-D\Delta+2R(2m-R)+\binom R2+2m
 &=-m-D\Delta+4mR-\frac{3R^2+R}{2}.
 \label{eq:residualexp}
\end{align}
The nondirect case has $R\ge1$. Every nonzero edge space has dimension two, so
$k\ge2$ and $R\le2m-2\le64$. By \eqref{eq:ceilbound},
\begin{equation}
 \Delta\ge j(R):=\left\lceil\frac{11R}{32}\right\rceil.
 \label{eq:j}
\end{equation}
Because the coefficient of $m$ in \eqref{eq:residualexp} is $4R-1>0$, replacing
$m$ by 33 gives a valid upper bound for the exponent:
\begin{equation}
 -Dj(R)+B(R),\qquad
 B(R)=\frac{263R-3R^2}{2}-33,
 \qquad 1\le R\le64.
 \label{eq:BR}
\end{equation}
This maximizes over a \emph{larger} collection than the realizable profiles.
Thus the resulting constants are safe sufficient bounds, not claims of
optimal geometric dimension thresholds.

\begin{corollary}[Improved nondirect residual estimate]
\label{cor:residual}
For the point graph $G_*$ and the residual gain interface above, the total
nondirect residual contribution satisfies
\begin{equation}
 \begin{cases}
 O_D(q^{-7}),&D=310,\\
 O_D(q^{-D+303}),&\text{every even }D\ge310,\\
 O_D(q^{-D+224}),&\text{every even }D\ge364.
 \end{cases}
 \label{eq:improvedresidual}
\end{equation}
The direct singular residual contribution is $O_D(q^{-m})$ on an $m$-edge
stratum, hence at most $O_D(q^{-1})$ after summing all singular direct residual
strata.
\end{corollary}

\begin{proof}
Appendix~\ref{app:exponents} gives an exact 22-row certificate, grouping all
$1\le R\le64$ by $j(R)$, for
\begin{equation}
 \max_R\{B(R)-310j(R)\}=-7,\qquad
 \max_R\{B(R)-364j(R)\}=-140.
 \label{eq:exponentmax}
\end{equation}
The maxima occur at $R=8$ and $R=2$, respectively. Since $j(R)\ge1$, increasing
$D$ from 310 or 364 decreases every upper bound at least linearly with slope
one. This gives both uniform formulas in \eqref{eq:improvedresidual}.
There are only finitely many edge-support and local-rank strata at fixed edge
dimensions, so summing preserves the estimates. For direct families use
Corollary~\ref{cor:weighted}: the exponent is $-3m+2m=-m$.
\end{proof}

The local common dimensions $c_j=0$ were not used in deriving
\eqref{eq:BR}; only the gain bound and $h_e\in\{0,2\}$ were used.
This intentional relaxation is one source of potential further improvements.

\subsection{What changes in the source proof, and what does not}

The source's nondirect residual estimate is $O_D(q^{-D+1641})$ and its final
singular-tail argument chooses an even $D$ above an earlier residue threshold
and additionally $D>1641$ \cite[\src{prop:singular-tail}]{OpenAISidorenko}.
With the present replacement, that \emph{additional residual restriction} can
be changed to $D\ge310$. The other dimension conditions remain in force.
For a fixed even $D\ge364$, the explicit exponent offset improves from 1641
to 224, a reduction of 1417.

This distinction is essential. The source has other sufficiently-large-dimension
requirements, including the reduction of large intersection profiles and the
transverse finite-field analysis. We have not bounded those thresholds by 310
or 364. Nor have we made the smallest admissible field size explicit, produced
a smaller 35-vertex pattern, or optimized the eventual finite host order.
The source's global $q$-rate can still be limited by nonresidual profiles, so a
$q^{-7}$ residual estimate is not a $q^{-7}$ estimate for the entire construction.

The structural gain is nevertheless more than a numerical adjustment. The old
argument pays a worst-case constant for every nondirect family. The new argument
retains the relation dimension $R$, the incidence deficit $\Delta$, and the
number of lost alternating-form constraints. It therefore extends to other
fixed graph patterns and prescribed edge dimensions.

\section{Verification, provenance, and integration}
\label{sec:verification}

\subsection{What has actually been checked}

The package contains a standard-library Python verifier, using exact integer,
rational, and prime-field arithmetic with seed 20261007. Its recorded run checks:
\begin{enumerate}[label=(\roman*)]
\item the 13-point rotation system, its 22 face cycles, all 8,100 vertex subsets
of size at least three, and the unique maximizing density $32/11$;
\item 1,200 subspace arrangements over $\mathbb F_2,\mathbb F_3,\mathbb F_5$,
including a directness-repair certificate for every arrangement;
\item 320 polymatroid--cover comparisons and 702 unit contractions, with
submodularity and monotonicity rechecked after every contraction;
\item all 15,750 assignments of subspaces of $\mathbb F_2^2$ to the edges of
$K_3$ and $K_4$, testing the equality classification;
\item 163 exterior-rank and exponent-cancellation checks, exact four-dimensional
symplectic counts over $\mathbb F_2$ and $\mathbb F_3$, and every residual
exponent in the finite optimization.
\end{enumerate}
All checks pass in the recorded run. The precise counts and grouped exponent
table are written to \src{certificates/verification_results.json}.
These tests are not used as substitutes for the general proofs. For example,
testing representable polymatroids alone would not justify
Theorem~\ref{thm:cover}; its nonrepresentable cases are covered by the integral
contraction proof and continuity argument.

No Lean executable was available in the working environment, and no new Lean
proof is claimed. The accompanying code is a verifier of finite certificates,
not a formalization of the full mathematics.

\subsection{Suggested integration boundary}

A suitable new location is
\begin{center}
\src{Combinatorics/Ramsey/Research/SharpIncidenceStability/}.
\end{center}
This is a proposed destination, not a statement that the folder already exists.
The article and exact tests should initially be integrated as a research result
with written proofs. The source Sidorenko paper and its Lean comparator should
remain unchanged until the replacement count has been reviewed and translated
into their actual interfaces.

A formalization can be split into four independent layers. First formalize
finite normalized integral submodular functions and the unit-contraction lemma.
Second formalize the cover increments and sharp constant, initially with
rational weights and a finite maximum. Third establish the linear quotient
identities, directness repair, and strict-balance classification. Fourth build
the exterior-power and finite symplectic counting layer, importing the first
three layers only through their precise numerical conclusions.

The source-specific application requires a small explicit graph certificate,
the exact residual gain hypothesis, and the elementary integer optimization.
It should not be proved by assuming the desired final vanishing statement.
The general covariance and entropy applications can be integrated independently
of the finite-field argument.

\subsection{Novelty and verification ledger}

The following distinctions should accompany any repository summary.
The broad fractional-subadditivity framework, zero-gap modularity, information-
theoretic stability principle, and random-Ramsey use of $m_2$ are prior work
\cite{MadimanTetali,EllisEtAl,Jakhar2025,Jakhar2026,RamseyTilings}.
The sharp support formula, the direct proof supplied here, the strict-balanced
linear equality classification, and the relation-sensitive count are the
specific results advanced in this memorandum. The comparison against the
reviewed sources has not located these exact statements there; this is not an
exhaustive priority determination across all literature.

The 35-vertex Sidorenko counterexample remains a result claimed by the cited
source, with the formalization scope described by its own record. Our main
inequalities do not depend on its correctness. Our improvement of its weighted
residual tail has an explicit imported interface and does not claim to verify
the full construction.

\section{Further research questions}
\label{sec:questions}

The exact-constant problem has been settled here under the stated hypotheses.
The following questions are subsequent problems rather than hidden assumptions
in the proofs.

\begin{question}[Intrinsic equality for non-strict patterns]
Can the flag certificate in Corollary~\ref{cor:flag} be replaced by an intrinsic,
order-free description of all equality arrangements when several supports
maximize $\mu_a$? A natural target is a decomposition organized by overlapping
density-maximizing subhypergraphs. The common-core theorem settles the
unique-maximizer case, but does not imply a simultaneous direct decomposition
for arbitrary overlaps.
\end{question}

A concrete first step is to classify graphs made by gluing two strictly
2-balanced graphs along an edge. Depending on their densities, the maximizers
can be one side, both sides separately, or their union. Exact examples in these
three regimes should distinguish genuine nesting from artifacts of a chosen
quotient flag.

\begin{question}[Stability of sharp equality, not merely directness]
For a strictly $a$-balanced hypergraph, let
$S=\mu_a\Delta-R$. Does small $S$ force a family close, under an explicitly
specified dimension-edit distance, to a common-core-plus-private arrangement?
Can the edit cost be bounded using $S$ and the positive density gap between the
whole hypergraph and its proper supports?
\end{question}

Theorem~\ref{thm:repair} does not answer this question: it controls the cost of
removing \emph{all} dependence by $R$, which may be large when a large common
core nearly attains sharp equality. A stability theorem should retain that
core rather than delete it. The flag proof controls bad steps, but turning
that information into simultaneous subspaces requires an additional argument.

\begin{question}[Entropy extremizers and bounded-conditioning determinants]
What are all distributions attaining the sharp coefficient
$\Kopt_\gamma$, beyond the common-random-variable examples? For positive-definite
matrices with uniformly bounded condition number, what is the optimal
replacement for $\Kopt_\gamma$ in \eqref{eq:detbound}?
\end{question}

The determinant sharpness construction approaches a singular matrix. It does
not settle a compact, bounded-condition-number problem. Local expansion around
independence gives a quadratic form on correlations, while the unrestricted
constant is controlled by rank-one common-information supports; comparing these
two mechanisms is a concrete direction.

\begin{question}[Geometry of the relation space]
Can the exponent $Rk+B_N-A_N$ in \eqref{eq:finecount} be bounded substantially
better on relation spaces realizable as pairwise intersections of Lagrangians?
Can one replace the full Grassmannian cost $Rk$ by the dimension of the relevant
relation variety?
\end{question}

The proof uses the entire kernel bound
$B_N\le Rk+\binom R2$ and ignores the favorable term $A_N$ in the explicit
application. Neither relaxation is expected to be universally tight in the
source's residual geometry. Exact small-field searches can suggest which
relation patterns are feasible, but a field-uniform dimension argument is
needed for a new asymptotic bound.

\begin{question}[Using the source's zero common-face intersections]
For $G_*$, can the additional conditions $c_j=0$ eliminate the profiles governing
our relaxed maxima, especially $(R,\Delta)=(2,1)$ and $(8,3)$? What is the actual
maximum of the weighted exponent over realizable residual profiles?
\end{question}

All triangles of $G_*$ are the 22 listed faces, as can also be checked from the
edge list. This makes a triangle-supported common line incompatible with a
zero common-face intersection. It does \emph{not} automatically eliminate a
triangle that appears only after quotienting: lifting quotient information is
the substantive obstacle. A proof resolving that obstacle could improve the
constants without changing the analytic rank-layer construction.

\begin{question}[The full dimension threshold]
Can the earlier large-profile reduction and transverse counting thresholds in
the Sidorenko construction be made explicit and brought into the same numerical
range as the residual threshold? Can the field-size dependence then be made
effective enough to produce a finite host certificate?
\end{question}

This requires tracking the source's local gain parameters, its dimension-
independent bounded-profile reduction, and its subsequent finite-field errors.
The present residual improvement supplies one component of such a program,
not its completion.

\begin{question}[Beyond pair-incidence geometry]
What is the correct analogue of the exterior count for hypergraph-indexed
subspaces or for other classical polar spaces? When two hyperedges share
several vertices, the wedge multiplicities in \eqref{eq:ellzero} change.
Can their exact correction be combined with $\mu_a$ to give comparably sharp
singular-tail bounds?
\end{question}

The abstract stability theorem already handles uniform hypergraphs. What is
missing is not the entropy or span estimate, but a replacement for the
simple-graph wedge identity and a suitable geometric containment model.

\begin{question}[Formal proofs and reusable certificates]
Can the sharp fractional-cover theorem and strict-balanced equality theorem be
formalized as reusable finite-combinatorial results, independently of a
particular analytic application? Can the relation-sensitive count be organized
so that the source graph enters only through a density certificate and local
multiplicities?
\end{question}

A successful implementation would make future graph changes auditable without
repeating the entire proof. The finite support formula and the explicit
network reduction also suggest certificates for optimal constants: a maximizing
support gives the lower certificate, and minimum-cut or density inequalities
give the upper certificate.

\appendix
\section{A planar rotation certificate for the source graph}
\label{app:rotation}

The following cyclic neighbor orders specify an orientable rotation system.
For a directed edge $(u,v)$, move to $(v,w)$ where $w$ is the successor of $u$
in the row for $v$. This permutation of directed edges has exactly 22 cycles,
each of length three, and their underlying unordered triples are
\eqref{eq:faces}. Thus the data are independently checkable without a graph
layout algorithm.

\begin{center}
\begin{tabular}{@{}rl@{}}
\toprule
Vertex & Cyclic order of neighbors\\
\midrule
0 & $1,9,2,3$\\
1 & $0,3,10,7,12,9$\\
2 & $3,0,9,4$\\
3 & $1,0,2,4,11,10$\\
4 & $2,9,8,11,3$\\
5 & $7,10,11,6$\\
6 & $11,8,12,7,5$\\
7 & $12,1,10,5,6$\\
8 & $9,12,6,11,4$\\
9 & $4,2,0,1,12,8$\\
10 & $5,7,1,3,11$\\
11 & $10,3,4,8,6,5$\\
12 & $8,9,1,7,6$\\
\bottomrule
\end{tabular}
\end{center}

The graph is connected. Gluing oriented disks along the face cycles yields a
connected closed orientable surface with Euler characteristic two, hence a
sphere. This proves planarity. The independent density check in the verifier
enumerates all $2^{13}$ vertex subsets, ignores those of size below three, and
checks $e(S)\le3|S|-6$ on the remaining 8,100 subsets. For fixed vertex support,
adding missing induced edges only increases $(e-1)/(v-2)$, so this is also a
complete certificate of the maximum over edge supports. Only the entire
13-vertex set attains $32/11$, and it must retain all 33 edges to do so.

\section{Exact certificate for the residual exponent optimization}
\label{app:exponents}

Recall $j(R)=\lceil11R/32\rceil$ and
$B(R)=(263R-3R^2)/2-33$, for $1\le R\le64$.
The table gives the maximum of $B$ on each consecutive interval with fixed
$j$. It can be checked from the quadratic formula alone; its first difference
is $B(R+1)-B(R)=130-3R$, which changes sign only once.
The last two columns certify \eqref{eq:exponentmax}.

\begin{center}
\small
\begin{tabular}{@{}rrrrr@{}}
\toprule
$j$ & Range of $R$ & $\max B(R)$ & $\max B(R)-310j$ & $\max B(R)-364j$\\
\midrule
1 & $1$--$2$ &224&$-86$&$-140$\\
2 & $3$--$5$ &587&$-33$&$-141$\\
3 & $6$--$8$ &923&$-7$&$-169$\\
4 & $9$--$11$ &1232&$-8$&$-224$\\
5 & $12$--$14$ &1514&$-36$&$-306$\\
6 & $15$--$17$ &1769&$-91$&$-415$\\
7 & $18$--$20$ &1997&$-173$&$-551$\\
8 & $21$--$23$ &2198&$-282$&$-714$\\
9 & $24$--$26$ &2372&$-418$&$-904$\\
10 & $27$--$29$ &2519&$-581$&$-1121$\\
11 & $30$--$32$ &2639&$-771$&$-1365$\\
12 & $33$--$34$ &2704&$-1016$&$-1664$\\
13 & $35$--$37$ &2779&$-1251$&$-1953$\\
14 & $38$--$40$ &2827&$-1513$&$-2269$\\
15 & $41$--$43$ &2848&$-1802$&$-2612$\\
16 & $44$--$46$ &2849&$-2111$&$-2975$\\
17 & $47$--$49$ &2834&$-2436$&$-3354$\\
18 & $50$--$52$ &2792&$-2788$&$-3760$\\
19 & $53$--$55$ &2723&$-3167$&$-4193$\\
20 & $56$--$58$ &2627&$-3573$&$-4653$\\
21 & $59$--$61$ &2504&$-4006$&$-5140$\\
22 & $62$--$64$ &2354&$-4466$&$-5654$\\
\bottomrule
\end{tabular}
\end{center}

At $D=308$, the relaxed bound has value zero at $R=11$. At $D=309$, its maximum
is $-4$, again at $R=11$. The source uses even $D$, so 310 is the first even
dimension at which this particular relaxed optimization proves decay.
This is a statement about the displayed upper-bound calculation, not a lower
bound on the dimension required by the actual geometry.

\section{Reproduction and source identifiers}
\label{app:reproduce}

From the package root, run
\begin{verbatim}
python3 code/verify.py
latexmk -pdf -interaction=nonstopmode -halt-on-error article.tex
\end{verbatim}
The Python program uses only the standard library. The LaTeX document uses
standard TeX Live packages, including \src{newpxtext} and \src{newpxmath}.
The PDF is provided so that compilation is not required to read the result.

The mathematical source locations most relevant to the application are under
\begin{center}
\small
\src{preprints/A-counterexample-to-Sidorenkos-conjecture-September-23-2026/build/}.
\end{center}
Their blob identifiers at the pinned OpenAI commit are:
\begin{center}
\small
\begin{tabular}{@{}ll@{}}
\toprule
Source & Git blob identifier\\
\midrule
\src{sections/complex.tex} & \src{178b0e9da21867a727ea847dae8e537b9d719a70}\\
\src{sections/kernel.tex} & \src{bbe11be8fffe968012060ec43e6aff597e922831}\\
\src{sections/tail.tex} & \src{dc11e07a592e70d3a13e3a005c72378c88edc351}\\
\bottomrule
\end{tabular}
\end{center}

The package does not redistribute the upstream manuscripts or their font files.
It includes the finite graph data needed to reproduce the calculations, with
attribution, and the new article and verifier. The source code's finite checks
must not be described as an axiom audit or a theorem-prover certificate.

\begingroup
\small
\raggedright
\begin{thebibliography}{20}

\bibitem{OpenAIReadme}
OpenAI, \emph{Mathematics research collection: repository README},
\src{openai/math}, snapshot October 6, 2026.
\url{https://github.com/openai/math/blob/adc7f1241b42e322a6451854ab7e4b4c146bf78a/README.md}.

\bibitem{OpenAICatalog}
OpenAI, \emph{OpenAI Research Catalog}, October 6, 2026,
\src{overview.tex}, especially combinatorics families 159--161 and 170.
\url{https://github.com/openai/math/blob/adc7f1241b42e322a6451854ab7e4b4c146bf78a/overview.tex}.

\bibitem{ProveItStatus}
Vladimir Reshetnikov and ProveIt contributors,
\emph{Gowers Szemeredi formalization checkpoint}, October 7, 2026,
\src{Combinatorics/Ramsey/FORMALIZATION_STATUS.txt}.
\url{https://github.com/VladimirReshetnikov/ProveIt/blob/1210e9d3ac88a573f98d7ef8626ba3d2c30e9749/Combinatorics/Ramsey/FORMALIZATION_STATUS.txt}.

\bibitem{OpenAISidorenko}
OpenAI, \emph{A counterexample to Sidorenko's conjecture},
OpenAI Math Release preprint, September 23, 2026.
Source snapshot at commit \oaicommit.
\url{https://github.com/openai/math/tree/adc7f1241b42e322a6451854ab7e4b4c146bf78a/preprints/A-counterexample-to-Sidorenkos-conjecture-September-23-2026}.

\bibitem{OpenAILean161}
OpenAI, \emph{Counterexamples to Sidorenko's conjecture and the forcing
conjecture: scope of the Lean formalization}, \src{lean/docs/161.md},
October 2026 snapshot.
\url{https://github.com/openai/math/blob/adc7f1241b42e322a6451854ab7e4b4c146bf78a/lean/docs/161.md}.

\bibitem{MadimanTetali}
M. Madiman and P. Tetali,
\emph{Information inequalities for joint distributions, with interpretations
and applications}, IEEE Transactions on Information Theory \textbf{56} (2010),
2699--2713. DOI: \href{https://doi.org/10.1109/TIT.2010.2046253}{10.1109/TIT.2010.2046253}.
\url{https://arxiv.org/abs/0901.0044}.

\bibitem{EllisEtAl}
D. Ellis, E. Friedgut, G. Kindler, and A. Yehudayoff,
\emph{Geometric stability via information theory},
Discrete Analysis \textbf{2016}:10, 28 pp.
DOI: \href{https://doi.org/10.19086/da.784}{10.19086/da.784}.
\url{https://arxiv.org/abs/1510.00258}.

\bibitem{Jakhar2025}
G. Jakhar, G. R. Kurri, S. Chillara, and V. M. Prabhakaran,
\emph{Fractional subadditivity of submodular functions: equality conditions and
their applications}, arXiv:2501.12058v1, January 21, 2025.
\url{https://arxiv.org/abs/2501.12058}.

\bibitem{Jakhar2026}
G. Jakhar, G. R. Kurri, S. Chillara, and V. M. Prabhakaran,
\emph{Generalized information inequalities via submodularity, and two
combinatorial problems}, arXiv:2601.15723v1, January 22, 2026.
\url{https://arxiv.org/abs/2601.15723}.

\bibitem{RamseyTilings}
L. Arag\~ao, X. Cheng, R. Filipe, R. Miyazaki, D. Peng, and Z. Yan,
\emph{Ramsey properties for tilings in random graphs},
arXiv:2605.21471v1, 2026; see Section 1.1 for the 2-density convention and
its random Ramsey context.
\url{https://arxiv.org/abs/2605.21471}.

\end{thebibliography}
\endgroup
\end{document}
